\documentclass[12pt,a4paper]{article}
\usepackage{amsmath,amssymb,amsthm}
\usepackage{enumitem}
\usepackage{hyperref}
\usepackage{geometry}
\geometry{margin=2.5cm}

\newtheorem{theorem}{Theorem}[section]
\newtheorem{corollary}[theorem]{Corollary}
\newtheorem{proposition}[theorem]{Proposition}
\theoremstyle{definition}
\newtheorem{definition}[theorem]{Definition}
\theoremstyle{remark}
\newtheorem{remark}[theorem]{Remark}

\title{Big Bang Initial Condition from Recognition Science:\\
Zero-Defect Ledger State and No-Singularity Resolution}

\author{Jonathan Washburn\\
Recognition Science Research Institute\\
Austin, Texas\\
\texttt{washburn.jonathan@gmail.com}}

\date{March 2026}

\begin{document}
\maketitle

\begin{abstract}
We derive the cosmological initial condition from Recognition Science (RS).
The origin of the universe is the unique zero-defect ledger configuration:
all entries equal unity, total defect vanishes, and this state is the global
minimum of the cost functional.
The Big Bang is interpreted as the first nonzero-defect activation (first tick),
not a physical singularity with divergent density.
We show that Penrose's low-entropy puzzle is resolved: the initial state is
the unique cost minimizer, so low entropy is a theorem rather than a hypothesis.
\end{abstract}

\section{Introduction}

\subsection{The Singularity Problem in Standard Cosmology}

Standard cosmology traces the expansion of the universe backward in time.
Under general relativity, the Penrose--Hawking singularity theorems
\cite{hawking1970singularities,penrose1979singularities} establish that
under generic conditions (reasonable energy conditions, causal structure),
a past-incomplete geodesic implies an initial singularity: a boundary
at which curvature and density diverge.
The Friedmann equations yield $\rho \propto a^{-3(1+w)}$, so as $a \to 0$,
density $\rho \to \infty$ and the classical description breaks down.

This raises a fundamental question: what happened at $t = 0$?
The singularity is not a point in spacetime---it is a boundary beyond which
the theory does not extend.
Quantum gravity proposals (loop quantum cosmology, string cosmology, etc.)
attempt to resolve or replace the singularity, but typically introduce
new degrees of freedom or pre-Big-Bang phases.

\subsection{Recognition Science Approach}

Recognition Science (RS) provides a different resolution.
Rather than modifying general relativity or adding pre-inflationary physics,
RS derives the initial condition from the structure of the cost functional
itself.
The universe is modeled as a ledger: a discrete structure whose entries
evolve according to a cost that penalizes deviation from unity.
The initial state is then the unique configuration that minimizes this cost---a
finite, well-defined state with no singularity.

\section{Recognition Science Framework}
\label{sec:framework}

All results follow from the unique cost functional determined by
three axioms~\cite{RS_Axioms}.

\noindent\textbf{Axiom A1 (Normalization):} $J(1) = 0$.\\
\textbf{Axiom A2 (Recognition Composition Law):}
$J(xy)+J(x/y)=2J(x)J(y)+2J(x)+2J(y)$ for all $x,y>0$.\\
\textbf{Axiom A3 (Calibration):} $J''_{\log}(0) = 1$.

\begin{theorem}[Cost Uniqueness]
\label{thm:cost-unique}
The continuous solution to \textup{(A1)--(A3)} is unique:
$J(x) = \tfrac{1}{2}(x + x^{-1}) - 1$.
\end{theorem}

\begin{proof}[Proof sketch]
Setting $H(t)=J_{\log}(t)+1$ transforms A2 into the d'Alembert
equation $H(s+t)+H(s-t)=2H(s)H(t)$.  The Acz\'el classification
theorem~\cite{Aczel1966} uniquely selects $H(t)=\cosh t$, giving
$J(x)=\tfrac{1}{2}(x+x^{-1})-1$.
\end{proof}

\noindent The global minimizer of $J$ is $x=1$ with $J(1)=0$.
This zero-cost ground state is the RS initial condition.

\section{Cost Functional and Total Defect}

\begin{definition}[Cost Functional]
For $x > 0$, the RS cost functional is
\begin{equation}
J(x) = \frac{1}{2}\left(x + x^{-1}\right) - 1.
\end{equation}
\end{definition}

\begin{definition}[Configuration and Total Defect]
A \emph{configuration} of size $N \geq 1$ is a tuple $c = (x_1, \ldots, x_N)$
with each $x_i > 0$.
The \emph{total defect} (or defect functional) is
\begin{equation}
D(c) = \sum_{i=1}^{N} J(x_i).
\end{equation}
\end{definition}

\begin{definition}[Unity Configuration]
The \emph{unity configuration} $c_{\mathrm{unity}}$ is the configuration
in which every entry equals 1: $x_i = 1$ for all $i$.
\end{definition}

\section{Initial-Condition Theorems}

\begin{theorem}[Zero Cost at Unity]
$J(1) = 0$, and $J(x) > 0$ for all $x > 0$ with $x \neq 1$.
\end{theorem}

\begin{proof}
Direct computation: $J(1) = \frac{1}{2}(1 + 1) - 1 = 0$.
For $x \neq 1$, the arithmetic--geometric inequality gives
$x + x^{-1} > 2$ (strict for $x \neq 1$), hence
$J(x) = \frac{1}{2}(x + x^{-1}) - 1 > 0$.
\end{proof}

\begin{corollary}[Zero Defect of Unity Configuration]
$D(c_{\mathrm{unity}}) = 0$.
\end{corollary}

\begin{proof}
Each entry of $c_{\mathrm{unity}}$ equals 1, so $J(x_i) = 0$ for all $i$.
Thus $D(c_{\mathrm{unity}}) = \sum_i J(1) = 0$.
\end{proof}

\begin{theorem}[Uniqueness of Zero-Defect State]
If $D(c) = 0$, then $c = c_{\mathrm{unity}}$.
\end{theorem}

\begin{proof}
$D(c) = \sum_{i=1}^{N} J(x_i)$ is a sum of non-negative terms ($J(x_i) \geq 0$).
Hence $D(c) = 0$ if and only if $J(x_i) = 0$ for every $i$.
By Theorem 1, $J(x_i) = 0$ if and only if $x_i = 1$.
Therefore every entry of $c$ equals 1, so $c = c_{\mathrm{unity}}$.
\end{proof}

\begin{theorem}[Global Minimum]
The unity configuration is the unique global minimizer of $D$.
\end{theorem}

\begin{proof}
$D(c) \geq 0$ for all configurations, and $D(c) = 0$ only for $c = c_{\mathrm{unity}}$
by Theorem 2.
Thus $c_{\mathrm{unity}}$ attains the minimum value 0 and is the only configuration
that does so.
\end{proof}

\begin{remark}
These results hold for any finite $N$.
The initial-condition structure is independent of ledger size.
\end{remark}

\section{No-Singularity Resolution}

The RS framework replaces the cosmological singularity with a well-defined
initial state:

\begin{enumerate}[label=(\roman*)]
\item \textbf{Initial total defect is zero.} The universe begins at $D = 0$.
\item \textbf{That state is globally minimal.} No configuration has lower cost.
\item \textbf{The state is unique.} There is no ambiguity about the initial condition.
\end{enumerate}

Therefore RS does not require an infinite-density boundary condition.
The origin is a unique, finite, minimum-cost ledger state.
There is no divergence of curvature or density---only the first transition
away from the zero-defect baseline.

\section{Penrose's Low-Entropy Puzzle}

Penrose \cite{penrose1979singularities,penrose1989emperor} emphasized a deep
puzzle: the universe began in an extraordinarily low-entropy state.
The phase-space volume of configurations that could have evolved into our
observed universe is fantastically small compared to the volume of
``typical'' high-entropy states.
Why did the universe start in such a special condition?

In standard cosmology, this is treated as a hypothesis or fine-tuning.
In RS, it is a \emph{theorem}.
The initial state is the unique minimizer of the cost functional $D$.
By construction, the zero-defect configuration has the lowest possible
entropy (identified with total defect in RS).
No statistical reasoning, no measure theory on phase space, and no
fine-tuning are required---the low-entropy initial condition follows from
the uniqueness of the cost minimizer.

\section{The First Tick: Physical Interpretation}

What does ``first tick'' mean physically?

In RS, time advances via \emph{recognition events}: discrete updates to the
ledger in which configurations are identified and the cost is reduced.
At $t = 0$, the ledger is in the unity configuration---no recognition events
have occurred, no causal structure has been established.

The \emph{first tick} is the first recognition event: the first moment at
which some ledger entry deviates from unity, activating a nonzero defect.
This is not a singularity---it is the first discrete step in a well-defined
dynamics.
The 8-tick epoch ($\tau = 2^D$ with $D = 3$) structures the subsequent
evolution, but the initial condition itself requires no prior physics.

\section{Predictions and Consequences}

The RS initial-condition framework yields several predictions:

\begin{enumerate}[label=(\roman*)]
\item \textbf{No pre-Big-Bang physics.} The zero-defect state is the absolute
starting point; there is no preceding phase to explain or model.
\item \textbf{No bouncing cosmology.} The framework does not require a
contracting phase before expansion; the first tick initiates expansion from
a finite baseline.
\item \textbf{No initial singularity.} Density and curvature remain finite;
the ``origin'' is a minimum-cost ledger state, not a divergent boundary.
\item \textbf{Deterministic initial condition.} The unity configuration is
unique; there is no ensemble of possible initial states to average over.
\end{enumerate}

These predictions are falsifiable: observation of a pre-Big-Bang phase,
a bounce, or evidence for an initial singularity would challenge the RS
picture.

\section{Scope and Limitations}

\begin{itemize}
\item \textbf{Proved:} Uniqueness and minimality of the zero-defect initial
state; $J(1) = 0$ and the implication $D(c) = 0 \Rightarrow c = c_{\mathrm{unity}}$.
\item \textbf{Interpreted:} The no-singularity reading of the cosmological
origin; the identification of the Big Bang with the first recognition event.
\item \textbf{Deferred:} Full observational reconstruction of the earliest
post-activation epochs (inflation, reheating, BBN) is treated in companion
papers on the early universe.
\end{itemize}

\section{Conclusion}

Recognition Science closes the cosmological initial-condition problem as a
zero-defect theorem.
The universe begins at a unique minimum-cost ledger state: all entries equal
unity, total defect zero.
The Penrose--Hawking singularity is replaced by a finite, well-defined
baseline; Penrose's low-entropy puzzle is resolved because the initial state
is the unique cost minimizer.
The Big Bang is the first tick---the first recognition event away from that
baseline---not a physical singularity with divergent density.

\begin{thebibliography}{9}
\bibitem{RS_Axioms}
J.~Washburn, ``The Algebra of Reality: A Recognition Science Derivation
of Physical Law,'' \emph{Axioms} \textbf{15}(2), 90 (2026).

\bibitem{Aczel1966}
J.~Acz\'el, \textit{Lectures on Functional Equations and Their Applications},
Academic Press, New York (1966).

\bibitem{hawking1970singularities}
S.\,W. Hawking and R. Penrose,
``The singularities of gravitational collapse and cosmology,''
\textit{Proc. R. Soc. Lond. A} \textbf{314}, 529--548 (1970).

\bibitem{penrose1979singularities}
R. Penrose,
``Singularities and time-asymmetry,''
in \textit{General Relativity: An Einstein Centenary Survey},
eds. S.\,W. Hawking and W. Israel (Cambridge University Press, 1979), pp. 581--638.

\bibitem{penrose1989emperor}
R. Penrose,
\textit{The Emperor's New Mind} (Oxford University Press, 1989).
\end{thebibliography}

\end{document}
